\section{Transductive generalisation and its limits}
\label{sec:generalisation}

\subsection{Three different spaces can change}

A transductive investigation involves at least three spaces: physically possible realisations, realisations reachable from the present supported state, and realisations currently justified by the available evidence. A new measurement may change the third without changing the first. An authorised construction may enlarge the second without changing physical law. A misinterpreted measurement may only change an assessor's belief.

LM3 explicitly separates the actual feasible set from its estimate and distinguishes recognition, determination and execution \cite[\S17.1]{LM3}. Apoautotic language should retain that separation. An organisation has not acquired an interface because its model names one, and it has not established a continuation merely because its model predicts that one exists.

The same distinction prevents a misleading all-or-nothing view of the faculty. A system might reliably classify a narrow family of already described candidates, yet fail to originate useful candidates outside that family. Another might propose valuable candidates but be unable to validate them. A third might validate supplied candidates while relying on external construction. Such partial capacities are research results at their own scope, not evidence of the whole faculty.

\subsection{Generalisation is conditional, not supplied by common notation}

Let a finite dataset assign outcomes on \(D\subsetneq X\), with the hypothesis class containing every Boolean map on \(X\). Any unobserved \(x_*\) admits two hypotheses that agree on \(D\) and differ at \(x_*\). Therefore no learner using only those observations can guarantee the right answer there for every hypothesis. LM3 uses this elementary counterexample to delimit an unrestricted generative representation claim \cite[Proposition 11.2]{LM3}.

It does not follow that useful discovery is impossible. The additional basis may be a justified restricted model class, invariances, measurements of the new candidate or well-supported physical constraints. The research burden is to state that basis and show that it survives the proposed use. A faculty that learns from heterogeneous examples need not assume they all instantiate one context-free ``language of matter.''

An empirical transfer claim must also separate a new instance from a near-duplicate account of an old one. Reports sharing an experimental root cannot become independent validations merely because they are written differently. The proposed representation is useful only if it preserves the distinctions needed for the held-out tasks, including failures and warranted abstention, at matched cost \cite[\S\S10.4,11.4,20.6]{LM3}.

\subsection{No universal decision oracle is implicit in the faculty}

LM3 proves that eventual distinguishability is undecidable for an unrestricted class of finitely described deterministic programmes with unbounded memory and unspecified finite horizon. Its reduction makes the input bit observable only after an arbitrary simulated computation halts. A terminating decision procedure for every such candidate would decide halting \cite[Theorem 17.2]{LM3}.

That theorem neither eliminates bounded finite-state assessment nor proves a physical claim about unlimited memory. It blocks a particular universal promise. More agents, a richer terminology or self-revision do not remove the logical obstruction while retaining the same unrestricted model class. An apoautotic faculty may be partial, bounded, probabilistic or equipped to abstain. Those properties must be reported rather than treated as an embarrassment to hide.

A capability-preservation theorem can therefore remain meaningful without promising universal success. The preserved faculty may guarantee safe assessment and honest refusal for a specified class, rather than guaranteed construction for every conceivable substrate. Its negative outputs can carry indispensable information: the candidate is not supported by present evidence, the route is outside the permitted repertoire, or a stated contract has a counterexample. Those are different outcomes.

\subsection{Discovery does not remove physical accounting}

The mathematical results in this paper assume descriptions of transitions and observations. A physical transductive faculty must obtain and maintain their correspondence to reality. Preparation, measurement, calibration, memory, control, energy, timing and upkeep are part of the assessment. The absence of one familiar host does not make their cost vanish \cite[\S\S6,13,16]{LM3}.

Nor should this burden be restated as a premature impossibility claim. Not possessing a certificate for an unfamiliar medium does not prove that no adequate certificate or implementation exists. The appropriate distinction is between a counterexample to a specified claim, an inability to establish it from present evidence, and a theorem excluding an entire declared model class. The same distinction is necessary on the positive side: an existence proof for a model is not a measured transducer.

\section{Generation, organisational continuity and identity}
\label{sec:identity}

\subsection{A successor need not be a copy}

The \(\mathsf G\) coordinate asks whether a successor organisation or realisation has been established. The \(\mathsf C\) coordinate asks what continuity relation it bears to the predecessor. A differently parameterised implementation can continue the same behavioural contract, while a bit-for-bit copy can fail to continue the actual execution if it loses pending inputs, calibration, timing or external effects. Conversely, a successor can be generated without preserving the predecessor's task.

In this paper, these observations motivate an operational separation, not a universal metaphysics of generation. The integration directive identifies \(\mathsf G\) as a distinct concern but does not supply a complete formal definition \cite{REC}. We therefore do not infer reproductive capacity, novelty of organisation or creation of a new subject from the generation of an executable candidate.

Apoautosis and reproduction can overlap, but neither term should erase the other. One authorised successor can replace a predecessor without multiplying active instances. A copying process can produce additional instances while none possesses an apoautotic faculty. A generated successor may itself retain a capacity to generate further successors. The last is an additional capability claim, not a semantic consequence of the word \emph{copy}.

\subsection{Branching exposes a limit of operational individuation}

The source dialogue presents a predecessor \(A\) producing two successors \(B\) and \(C\). Both receive the consequential state, satisfy the same contract and retain the machinery needed for further continuation. They later receive different inputs. Those facts can establish two organisational continuations. They do not establish that both are numerically the same individual as \(A\) \cite[07:18:51--07:22:32]{D}.

The point can be stated with minimal logic. If identity is ordinary numerical identity and \(A=B\) and \(A=C\), transitivity gives \(B=C\). If the successors are distinct, operational continuity cannot be identified with both asserted numerical identities. A branching continuation relation need not be an equivalence relation and is not thereby defective. It is answering a different question.

Removing branching does not settle the issue automatically. A unique successor can be the unique appropriate continuation under a policy without a theorem that organisational preservation constitutes personal identity. Additional metaphysical premises would be needed for that interpretation. The companion accordingly claims contract-relative continuity and, where supplied, lineage; it does not derive consciousness, phenomenal acquaintance or an indivisible bearer from a simulation relation.

\subsection{Three boundaries that remain different}

Historical origin identifies where a process began. Present support identifies what currently realises and sustains it. An individuation criterion states what counts as the continuing object for the inquiry. The user explicitly separates these in the source dialogue \cite[06:57:37]{D}. A process can originate on \(H_0\), continue elsewhere and still possess a historical relation to \(H_0\) without requiring its ongoing support.

Likewise, operational authority may select an authoritative continuation without deciding metaphysical identity. A version identifier can distinguish records without individuating a conscious subject. A conserved physical inventory can survive changes of constituent kinds without preserving every molecular object. None of these observations licenses a slide from ``some relation persists'' to ``everything constitutive of identity persists.''

The programme's metaphysical inquiries remain relevant but separate. A technical account can pose the bearer question sharply; it cannot answer it solely by accumulating more operational continuity evidence. Source-relative philosophical arguments must be identified as such, and their conclusions cannot be imported into a common-premise formal result without the additional premises being made explicit.

\section{Related work and the exact place of the synthesis}
\label{sec:related}

The mathematical methods used here have substantial ancestry. The paper does not claim that a new name makes refinement, sufficient information, state identification, safety games or organisational closure newly discovered. The distinction to assess is the \emph{combined target}: preservation or lawful reacquisition of a declared future transductive faculty while the organisation's present implementation and required supports change.

The following is a bounded related-work check, not a comprehensive priority survey. Published abstracts or institutional records were checked for the named formal-methods sources; the author-hosted text of Mont\'evil and Mossio was also consulted. No empirical study is replicated, and the present paper's proofs do not depend on an unexamined assertion of a new theorem in that literature.

\subsection{Refinement and contract-based design}

Abadi and Lamport study mappings from lower-level to higher-level specifications, including the relation between state-machine steps and permitted behaviours. Their work is a direct antecedent for the commuting-map style of continuation used here \cite{AL}. LM3 already identified this ancestry. Our finite induction is not a replacement for their theory, and it supplies no general completeness theorem for arbitrary apoautotic systems.

Benveniste and colleagues provide an abstract contract framework supporting verification, abstraction, refinement and synthesis across system-design methods \cite{CONTRACT}. The presence of assumptions, guarantees and composition conditions in the current paper therefore cannot itself establish originality. A claimed advance would need to show a result or practical benefit not already obtained from an appropriately instantiated contract account.

The particular question added here is whether a contract for \emph{this service} protects the observations and realisation capacities needed for \emph{later continuation tasks}. The calibration example shows that it need not. This is a choice of explanandum and composition burden, not a claim that ordinary contracts are unable to express it.

\subsection{Adaptive software and state identification}

Vogel and Giese's EUREMA work concerns engineering adaptation engines and feedback loops for self-adaptive software \cite{EUREMA}. That establishes a clear neighbouring research area for changing the machinery that changes a system. It does not, merely by being an adaptation architecture, establish the broader physical-interface or cross-support claims discussed here.

Van den Bos and Vaandrager generalise adaptive state-identification techniques to labelled transition systems with inputs and outputs, addressing compatible states and distinguishing tests \cite{ID}. The use of observation classes, adaptive inquiry and backward reachability in this paper must be compared with that tradition. Theorems~\ref{thm:belief} and~\ref{thm:requests} use standard finite constructions; their contribution here is to expose the original-target/current-executability and future-request obligations explicitly, not to claim a new general identification algorithm.

\subsection{Autopoiesis and closure of constraints}

Varela, Maturana and Uribe characterise a class of autopoietic systems and present a simulated minimal model \cite{AUTO}. The existence of that tradition prevents treating all organisation-maintaining processes as a new category invented here. Apoautosis is not proposed as a substitute definition of life.

Mont\'evil and Mossio distinguish processes from constraints conserved at relevant timescales, and describe organisational closure through dependencies among constraints. They explicitly keep such closure compatible with thermodynamic openness and environmental conditions \cite{CLOSURE}. That is a substantive conceptual neighbour of the present distinction between an organisation's constitutive technical roles and its ongoing causal support.

The question pursued here is more specifically about re-realisation and retained future capacity. An organisational account can describe maintenance in one support regime without providing a route into another or preserving the means for later transitions. Conversely, a technical continuation satisfying this paper's contract need not meet a biological theory's conditions for living organisation. Whether a stronger formal correspondence between these accounts can be established is open, not settled by the word \emph{closure}.

\subsection{Inherited versus newly assembled content}

LM3 supplies the physical-role decomposition, realisation/migration equations, joint-witness obligations, resource boundaries and maintained-operation distinction. The dialogue isolates the endogenous faculty and its retention. S2 supplies the information-preservation and reacquisition analysis; S3 consolidates its task-relative interpretation. GF1 supplies future-safe representation and landscape preservation; ICS5 supplies observation congruence and destructive inquiry; CDR supplies the fixed-target priced-probe account \cite{LM3,D,S2,S3,GF1,ICS,CDR}.

The current paper's contribution is to make these pieces answer one explicit question rather than to treat them as interchangeable results. Its finite-request composition theorem is a newly assembled exposition using conventional mechanisms. Its source-status reconciliation is also explicit: a preliminary dialogue, a partial consolidation and a conditional paper do not become independently validated facts by citing one another.

\section{A bounded research programme for apoautotic faculties}
\label{sec:research}

\subsection{The immediate mathematical destination}

The focused problem is this:

\begin{quote}
Given a present operational contract, a family of future realisation requests, restricted observations and interventions, and specified support changes, characterise which continuations preserve both warranted decision and executable response---and determine the least retained information, reacquisition effort and resource margin needed to do so.
\end{quote}

This is not equivalent to keeping all bits, identifying the complete world or finding a feasible endpoint. It joins four demands: retain the distinctions that matter, retain or obtain access to them, avoid destroying what subsequent decisions require, and reach a successor in which the next required operation remains possible and permitted.

The finite results here settle deliberately restricted pieces. Proposition~\ref{prop:cover} gives exact unrestricted record cardinality for one relation-valued decision. Proposition~\ref{prop:budget} gives a necessary fixed-target response budget. Theorem~\ref{thm:belief} characterises bounded decision under specified state-changing inquiry. Theorem~\ref{thm:requests} composes such tasks over a finite request horizon. None gives a general optimal resource trade-off for changing physical realisations with uncertain interfaces and accumulating costs.

A stronger result could identify an exact capability-preserving quotient, a sharp lower bound under restricted reacquisition, or a minimal condition for transferring a policy across changed state spaces. Such a result would need precise hypotheses, comparison with existing abstraction and game-solving methods, and a proof that its conclusion is not merely restating those hypotheses. The aim is not an ever-larger conjunction called ``apoautosis.''

\subsection{The immediate empirical destination}

The empirical sequence should separate service continuity, retained future capability and attributed endogeneity. A first bounded experiment could use two predeclared, authorised implementations of a small stateful task, with a future calibration or diagnostic request fixed in advance. It would compare three cases: preserve the required future record, discard it but retain a genuine reacquisition interface, and discard it without such an interface. The reported outcome would concern that task family and support boundary only.

The decisive negative controls follow directly from the paper. A retained display with lost consequential state should fail. A successful task with a lost future calibration should count as present success and future failure, not as total success or total failure. A reference that still depends on the removed host should not count as independent reacquisition. A correct readout with unavailable action resources should fail an executable-continuation claim. Permission withdrawal should produce the designated permitted termination, not an attempt to maintain activity regardless.

A separate endogeneity assessment would test which declared internal functions actually contribute to the result and which are supplied by an external technical operator. This is not established by replacing the operator's name in a diagram. The attribution must be supported by observable interfaces and appropriately controlled interventions. The test need not remove external legitimate oversight to investigate technical participation.

\subsection{Baseline and ablation requirements}

A fair evaluation should compare an ordinary-language contract account, a conventional formal model supplied with the same evidence, and the proposed apoautotic presentation. None should receive hidden extra information, a weaker target or an artificially poor implementation. Outcomes should distinguish false admission, false rejection, justified unresolved status, missed dependency, lost future capability and resource cost. Total cost includes building and maintaining the representation, not only the solver's query count.

With identical decision functions and inputs, changing terminology cannot change the verdict. LM3 already includes such equality as a control, not as evidence of superiority \cite[\S\S12.5,20.10]{LM3}. Any demonstrated advantage would instead concern, for example, more reliable specification of future obligations, earlier detection of lost access or less costly assessment under matched information. A no-difference outcome is informative and must not be relabelled a success.

The accompanying finite checker tests expressibility and internal consistency of the examples, not comparative human or model performance. Its counterexamples could be independently expressed in standard formal methods. The research claim becomes stronger only when the combined treatment produces an exact result, a useful bound, a sharper rejection or a measured advantage relative to the relevant baseline.

\subsection{The place of meniscus in this companion}

Apoautosis is one research front within Orthemology. It does not replace the programme's epistemic, metaphysical, linguistic or foundational fronts. Nor does a permitted cross-reference to OWOU, revision certificates or the CV calculi establish an equivalence. The supplied integration instruction explicitly warns against those unproved identifications \cite{REC}.

For this front, a potential meniscus-level result would be an independently significant characterisation or limitation of future-capability preservation across changing supports, or a validated construction within a nontrivial declared class. A broadly useful negative result could qualify as well as a positive one. The status would require serious prior-art comparison and independent specialist scrutiny; a research workflow cannot award it to itself because its internal checklist is complete.

No such break is claimed here. The paper provides a more definite object for that effort: it identifies the preserved faculty, gives finite conditions that can be tested, and exposes where current continuation evidence stops. A substantial future result would have to add to that object rather than merely attach a stronger label to the same examples.

\section{Evidence, authority and disclosure}
\label{sec:evidence}

The paper distinguishes conceptual proposal, ordinary conditional proof, executable finite check, empirical implementation and independent review. The supplied stage-two and stage-three reports retain incomplete overall-review statuses; those do not erase their stated local arguments, but they prevent presenting the reports as comprehensively certified integrations. The present manuscript is a new synthesis and not a retroactive approval of those reports \cite{S2,S3}.

Authorised research and successful implementation also remain separate. The mathematics here concerns specified state maps, information bounds and permitted transitions. It is not a live-system discovery engine, a covert-channel synthesis method, an autonomous resource-acquisition system or a mechanism for preserving a model against its operator's instructions. None of the finite checks changes an external system or searches for physical interfaces.

Constructive scientific results can have different disclosure implications from purely descriptive ones, but that possibility is not evidence of a present hazardous capability. No current hazard or emergency is established by the existence of these definitions. Any later result that materially enables consequential operations should have its implementation and disclosure reviewed separately from its mathematical correctness. Conversely, calling a claim sensitive must not shield its assumptions from appropriate scrutiny.

The operational boundary is therefore not ``the system promises to stop.'' Permission belongs to the external conditions under which the system is evaluated. A finite policy theorem assumes that boundary; it does not prove an implemented sandbox, control channel or revocation mechanism secure. Those are separate engineering and assurance obligations. This paper neither supplies nor certifies them.

\section{Conclusion}
\label{sec:conclusion}

The principal claim is not that an organisation can abandon material support. It is that its continuity need not be individuated by one particular support, while the faculty enabling further continuation is itself an object that can be preserved, transformed, lost or adequately reacquired. This turns a question about portable information into a question about supported organisational capacity.

LM3 supplied the distinction between a coherent description and an adequate, reachable, maintained realisation. The companion adds the successor question: after that boundary is crossed, what remains available for the next one? Preserving today's task does not answer it. The calibration example gives a finite countermodel. The record and query results identify which information claims are justified, while the belief-state and finite-request constructions explain why a successful local decision may still strand the future faculty.

The result is neither a universal construction protocol nor a new theory of personal identity. It is a disciplined account of the phenomenon and a set of conditional mathematical obligations. Transduction does not imply continuity; generation does not imply identity; a present result does not imply future capability; and a successful self-application does not supply its own warrant or authority.

The research direction is correspondingly specific: derive sharper conditions under which an organisation can continue through changing realisations while retaining the declared means to assess and establish further continuations. A future advance must supply the missing mathematical or physical evidence, not merely announce that the process has become recursive. Apoautosis names that richer preservation problem; apoautotic transductive faculties name the capacities whose persistence makes it more than a one-time transfer.
